\begin{exo}
	Une professeure de mathématiques a corrigé son paquet de $50$ copies du bac notées sur $20$.

	Le centre d'examen l'informe que la moyenne des correcteurs est de $10/20$.

	Elle doit vérifier si son paquet est plutôt meilleur ou moins bon que la moyenne du centre.

	Elle a enregistré ses notes dans un tableau.

	\begin{center}
		\def\arraystretch{1.5}
		\small
		\begin{tabular}{|l|c|c|c|c|c|c|c|}
			\hline
			\textbf{Valeur}   & 3  & 4  & 5  & 6  & 7  & 9  & 10 \\ \hline
			\textbf{Effectif} & 3  & 2  & 5  & 3  & 7  & 8  & 3  \\ \hline\hline
			\textbf{Valeur}   & 11 & 12 & 13 & 15 & 16 & 18 & 19 \\ \hline
			\textbf{Effectif} & 2  & 6  & 1  & 4  & 3  & 2  & 1  \\ \hline
		\end{tabular}
	\end{center}

	\begin{enumerate}
		\item Entrer les notes dans la calculatrice avec les effectifs de copies associés.
		\item Afficher sur la calculatrice l'ensemble des indicateurs statistiques puis donner la moyenne et l'écart-type. Où se situe son paquet de copies en comparaison à ceux du centre d'examen ? Arrondir au centième.
	\end{enumerate}
	
	Elle reprend ses copies et s'aperçoit qu'elle a enregistré un $9$ alors que le candidat avait $19$.

	\begin{enumerate}[resume]
		\item Modifier les deux effectifs concernés puis afficher de nouveau la moyenne. Au final, que peut-elle penser de son paquet ?
	\end{enumerate}
\end{exo}
